% ID: SEL-D0029-OA-E01
% ESTADO: validada. S030, cotejada con examen y criterios.
\textbf{a)} La matriz de la derecha elevada al cuadrado es $I_2$:
el elemento superior derecho es $1-1=0$ y la diagonal vale $1$.
Por tanto, con $X=(u,v)^t$, queda
$2u+3v=4$ y $u-5v=1$. Sustituyendo $u=1+5v$,
se obtiene $13v=2$. Así,
\[
\boxed{X=\begin{pmatrix}23/13\\2/13\end{pmatrix}}.
\]
Al multiplicar por la matriz de la izquierda se recupera $(4,1)^t$.

\textbf{b)} La resta $2A-3B$ exige que $B$ sea $3\times2$.
Como $AA^t$ es $3\times3$, $C$ debe ser cuadrada de orden 3
para que exista $C^2$ y se pueda restar. Para que exista $AD$,
$D$ debe tener dos filas. Puede tener cualquier número positivo
de columnas $k$, y el resultado será $3\times k$.
